\section{Anatomia do Dockerfile}

\begin{frame}
  \frametitle{O que é um Dockerfile?}
  Documento de texto com instruções para criar uma imagem de container.

  \vspace{0.5em}
  Especifica:
  \begin{itemize}
    \item Comandos a executar
    \item Arquivos a copiar
    \item Comando de inicialização
  \end{itemize}
\end{frame}

\begin{frame}
  \frametitle{Principais Instruções}
  \begin{itemize}
    \item \texttt{FROM}: Imagem base do build
    \item \texttt{WORKDIR}: Diretório de trabalho dentro da imagem
    \item \texttt{COPY}: Copia arquivos do \textit{host} para a imagem
    \item \texttt{RUN}: Executa um comando durante o build
    \item \texttt{ENV}: Define variável de ambiente no container
    \item \texttt{EXPOSE}: Indica a porta que a imagem deseja expor
    \item \texttt{USER}: Define o usuário padrão para comandos subsequentes
    \item \texttt{CMD}: Comando padrão ao iniciar o container
  \end{itemize}
\end{frame}

\begin{frame}
  \frametitle{Sequência Típica de um Dockerfile}
  \begin{enumerate}
    \item \texttt{FROM}: Determinar a imagem base
    \item \texttt{RUN}: Instalar dependências da aplicação
    \item \texttt{COPY}: Copiar código-fonte e binários
    \item \texttt{EXPOSE}, \texttt{USER}, \texttt{CMD}: Configurar a imagem final
  \end{enumerate}
\end{frame}

\begin{frame}
  \frametitle{Exemplo: Dockerfile Python}
  \texttt{FROM python:3.13}

  \texttt{WORKDIR /usr/local/app}

  \texttt{COPY requirements.txt ./}

  \texttt{RUN pip install --no-cache-dir -r requirements.txt}

  \texttt{COPY src ./src}

  \texttt{EXPOSE 8080}

  \texttt{RUN useradd app}

  \texttt{USER app}

  \texttt{CMD ["uvicorn", "app.main:app", "--host", "0.0.0.0", "--port", "8080"]}
\end{frame}

\begin{frame}
  \frametitle{Exemplo Prático: Aplicação Node.js}
  \texttt{FROM node:22-alpine}

  \texttt{WORKDIR /app}

  \texttt{COPY . .}

  \texttt{RUN yarn install --production}

  \texttt{CMD ["node", "./src/index.js"]}

  \vspace{0.5em}
  \begin{itemize}
    \item \texttt{node:22-alpine}: Node.js 22 com Alpine Linux (leve e enxuto)
    \item \texttt{yarn install --production}: Instala apenas dependências de produção
  \end{itemize}

  \vspace{0.5em}
  \begin{alertblock}{Atenção}
    O arquivo \texttt{Dockerfile} não tem extensão. Use \texttt{docker init}
    para gerar automaticamente um Dockerfile otimizado.
  \end{alertblock}
\end{frame}